\newcommand{\hmwkTitle}{Homework 2}
\newcommand{\hmwkDueDate}{Monday, January 11, 2016}
\newcommand{\hmwkHeaderTitle}{ICCS206 Discrete Mathematic: Homework 2}
\documentclass{article}

\usepackage{fancyhdr}
\usepackage{amsmath}
\allowdisplaybreaks
\usepackage{amsthm}
\usepackage{amsfonts}
\usepackage{tikz}
\usepackage[plain]{algorithm}
\usepackage{algpseudocode}
\usepackage{enumitem}
\usepackage{graphicx}
\usepackage{hyperref}

\usetikzlibrary{automata,positioning}

%
% Basic Document Settings
%

\topmargin=-0.45in
\evensidemargin=0in
\oddsidemargin=0in
\textwidth=6.5in
\textheight=9.0in
\headsep=0.25in

\linespread{1.1}

\pagestyle{fancy}
\lhead{\hmwkAuthorName}
\chead{} % empty
\rhead{\hmwkHeaderTitle}
\cfoot{\thepage}

\renewcommand\headrulewidth{0.4pt}
\renewcommand\footrulewidth{0pt}

\setlength\parindent{0pt}

%
% Homework Details
%

\providecommand{\hmwkTitle}{Homework}
\providecommand{\hmwkDueDate}{TBD}
\providecommand{\hmwkClass}{ICCS206 Discrete Mathematic}
\providecommand{\hmwkClassTime}{Section\ 1}
\providecommand{\hmwkClassInstructor}{Pakawut Jiradilok}
\providecommand{\hmwkAuthorName}{\textbf{Jiraroj Wiruchpongsanon (6781617)}}
\providecommand{\hmwkHeaderTitle}{\hmwkClass:\ \hmwkTitle}

%
% Title Page
%

\title{
    \vspace{2in}
    \textmd{\textbf{\hmwkClass:\ \hmwkTitle}}\\
    \normalsize\vspace{0.1in}\small{Due\ on\ \hmwkDueDate}\\
    \vspace{0.1in}\large{\textit{\hmwkClassInstructor\ \hmwkClassTime}}
    \vspace{3in}
}

\author{\hmwkAuthorName}
\date{}

\renewcommand{\part}[1]{\textbf{\large Part \Alph{partCounter}}\stepcounter{partCounter}\\}

%
% Helper Commands
%

\newcounter{partCounter}
\newcommand{\solution}{\textbf{\large Solution}}

\newtheorem*{theorem}{Theorem}

\theoremstyle{remark}
\newtheorem*{remark}{Remark}

\begin{document}

\maketitle
\newpage


\section*{Problem 1}
Use induction to prove
\[
1^2 + 2^2 + 3^2 + \cdots + n^2 = \frac{n(n+1)(2n+1)}{6}\qquad \forall n \in I^+.
\]
\medskip
\solution

 \begin{proof}

    we claim that
    \[
    P(n) = 1^2 + 2^2 + 3^2 + \cdots + n^2 = \frac{n(n+1)(2n+1)}{6}
    \]
    for the \textbf{base case} where $n_0 = 1$
    \[
    \begin{aligned}
    P(1) &= 1^2 = \frac{1(1+1)(2\cdot 1 + 1)}{6} = 1 \\
    P(2) &= 1^2 + 2^2 = 5 = \frac{2(2+1)(2\cdot 2 + 1)}{6} \\
    P(3) &= 1^2 + 2^2 + 3^2 = 14 = \frac{3(3+1)(2\cdot 3 + 1)}{6} 
    \end{aligned}
    \]
    then, we assume P(k) for arbitrary $(k \ge n_0)$, we prove $p(k+1)
  = \frac{(k+1)((k+1) + 1)(2(k+1) + 1)}{6}$
    \begin{align*}
    p(k+1) &= \underbrace{(1^2 + 2^2 + 3^2 + \cdots + k^2)}_{P(k)} +
  (k+1)^2\\
    &= \frac{k(k+1)(2k+1)}{6} + (k+1)^2 && \text{inductive hypothesis}\\
    &= \frac{k(k+1)(2k+1) + 6(k+1)^2}{6}\\
    &= \frac{(k+1)(k+2)(2k+3)}{6} && \text{factored the numerator}\\
    &= \frac{(k+1)((k+1) + 1)(2(k+1) + 1)}{6} && \text{rewrite in the expected form for } P(k+1)
    \end{align*}
    since P(k+1) follows P(k), the \textbf{induction case} is satisfied, thus we can conclude the given relation is true \end{proof}

\vspace{1em}
\hrulefill

\section*{Problem 2}
Use induction to prove
\[
1^3 + 2^3 + 3^3 + \cdots + n^3 = (1 + 2 + 3 + \cdots + n)^2, \quad \forall n \in I^+.
\]
\medskip
\solution

 \begin{proof}
  
    we claim that
    \[
        P(n) = 1^3 + 2^3 + 3^3 + \cdots + n^3 = (1 + 2 + 3 + ... + n)^2
    \]
    for the \textbf{base case} where $n_0 = 1$
    \[
    \begin{aligned}
    P(1) &= 1^3 = 1^2 = 1 \\
    P(2) &= 1^3 + 2^3 = 9 = (1 + 2)^2 \\
        P(3) &= 1^3 + 2^3 + 3^3 = 36 = (1 + 2 + 3)^2
    \end{aligned}
    \]
     then, we assume P(k) for arbitrary $(k \ge n_0)$, we prove $p(k+1) = (1 + 2 + 3 + ... + k + (k + 1))^2$
    \begin{align*}
    p(k+1) &= \underbrace{(1^3 + 2^3 + 3^3 + \cdots + k^3)}_{P(k)} + (k+1)^3 \\
        &= (1 + 2 + 3 + ... + k)^2+ (k+1)^3 && \text{inductive hypothesis} \\
        &= (\frac{k(k + 1)}{2})^2 + (k + 1)^3 && \text{well known arithmatic series identity --- } \sum_{i = 0}^n i = \frac{n(n + 1)}{2} \\
        &= (\frac{k(k + 1)}{2})^2 + (k^3 + 3k^2 + 3k + 1) && \text{binomial expansion} \\
        &= \frac{k^4 + 2k^3 + k^2 + 4(k^3 + 3k^2 + 3k + 1)}{4} \\
        &= \frac{k^4 + 6k^3 + 13k^2 + 12k + 4}{4} \\
        &= \frac{(k^2 + 3k + 2)^2}{4} && \text{factor the quartic} \\
        &= \left(\frac{(k+1)(k+2)}{2}\right)^2 && \text{factor the quadratic} \\
        &= \left(\frac{(k+1)((k+1) + 1)}{2}\right)^2 && \text{rewrite as } P(k+1)
    \end{align*}

    since P(k+1) follows P(k), the \textbf{induction case} is satisfied, thus we can conclude the given relation is true \end{proof}

 \end{proof}

\vspace{1em}
\hrulefill

\section*{Problem 3}
The Fibonacci numbers are defined by $F_0 = 1$, $F_1 = 1$, $F_2 = 2$, $F_{n+1} = F_n + F_{n-1}$. Prove that
\[
F_2 + F_4 + F_6 + \cdots + F_{2n} = F_{2n+1} - 1.
\]
\medskip
\solution

  \remark
  $0,1,1,2,3,5,8,13,21,34,55,89,144,233,377,610,987,1597,2584,4181$ ---
  few first fibonacci sequence from the OEIS

   \begin{proof}
    we claim that
    \[
        P(n) = F_2 + F_4 + F_6 + ... + F_{2n} = F_{2n+1} - 1.
    \]
    for the \textbf{base case} where $n_0 = 1$
    \[
    \begin{aligned}
    P(1) &= F_2 = 2 = F_{2(1)+1} - 1 = 3 - 1 \\
    P(2) &= F_2 + F_4 =  2 + 5 = F_{2(2)+1} - 1 = 8 - 1 = 7 \\
        P(3) &= F_2 + F_4 + F_6 = 2 + 5 + 13 = F_{2(3)+1} - 1 = 21 - 1
  = 20
    \end{aligned}
    \]
    then, we assume P(k) for arbitrary $(k \ge n_0)$, we prove $p(k+1)
  = F_{2(k+1)+1} - 1$
    \begin{align*}
    P(k+1) &= \underbrace{(F_2 + F_4 + F_6 + \cdots + F_{2k})}_{P(k)} +
  F_{2(k+1)}\\
        &= (F_{2k+1} - 1) + F_{2k+2} && \text{inductive hypothesis}\\
        &= (F_{2k+1} + F_{2k+2}) - 1 && \text{commutativity and
  associativity of addition in }\mathbb{R}\\
        &= F_{2k+3} - 1 && \text{fibonacci definition --- }F_{n+1} = F_n +
  F_{n-1}\\
        &= F_{2(k+1) + 1} - 1
    \end{align*}
    since P(k+1) follows P(k), the \textbf{induction case} is
  satisfied, thus we can conclude the given relation is true
   \end{proof}

  \vspace{1em}
  \hrulefill

\vspace{1em}
\hrulefill

\section*{Problem 4}
Recall that $a$ is divisible by $b$ if $\exists n \in I$ such that $a = nb$ (denoted $b \mid a$).
\begin{enumerate}[label=\arabic*.]
    \item Show that $8^n - 3^n$ is divisible by $5$ for all $n \in I^+$.
    \item For $a,b \in I$ and $n \in I^+$, use induction to show $a-b \mid a^n - b^n$.
    \item Use the previous result to show that if $n$ is odd then $a+b \mid a^n + b^n$.
\end{enumerate}
\medskip
\solution

 \begin{proof}
  
 \end{proof}

\vspace{1em}
\hrulefill

\section*{Problem 5}
Show that we can cover a $2^n \times 3^m$ checkerboard with L-shaped trominoes. (Hint: induct separately on $m$ and $n$.)
\medskip
\solution

 \begin{proof}
  
 \end{proof}

\vspace{1em}
\hrulefill

\section*{Problem 6}
Use induction to show that $n^3 \le 3^n$ for all $n \ge 1$. (Hint: prove for $n \ge 3$ first, then check $n=1,2$.)
\medskip
\solution

 \begin{proof}
  
 \end{proof}

\vspace{1em}
\hrulefill

\section*{Problem 7}
Use induction to show that
\[
\sum_{i=1}^n i \cdot 2^i = (n-1)2^{n+1} + 2.
\]
\medskip
\solution

 \begin{proof}
  
 \end{proof}

\vspace{1em}
\hrulefill

\section*{Problem 8}
Suppose you have an infinite supply of socks of five colors (green, red, black, yellow, white). You draw socks blindly one at a time. How many must you draw to guarantee at least one matching pair?
\medskip
\solution

 \begin{proof}
  
 \end{proof}

\vspace{1em}
\hrulefill

\section*{Problem 9}
Let $S$ be a set of $n+1$ integers. Show there exist $a,b \in S$ such that $a-b$ is a multiple of $n$. (Hint: work modulo $n$.)
\medskip
\solution

 \begin{proof}
  
 \end{proof}

\vspace{1em}
\hrulefill

\section*{Problem 10}
Prove that among any five points in a $1\text{m} \times 1\text{m}$ square, there is at least one pair whose distance is at most $\sqrt{2}/2$ meters. (Hint: partition the square into four equal squares.)
\medskip
\solution

 \begin{proof}
  
 \end{proof}

\vspace{1em}
\hrulefill

\end{document}
